\documentclass[11pt]{article}
\usepackage[margin=1in]{geometry}
\usepackage{amsmath,amssymb,mathtools}
\usepackage{fontspec,unicode-math}
\setmainfont{Latin Modern Roman}
\setsansfont{Latin Modern Sans}
\setmathfont{Latin Modern Math}
\usepackage{xurl,hyperref}
\hypersetup{colorlinks=true,urlcolor=blue,linkcolor=blue}
\newcommand{\sref}[2]{\hyperlink{source:#1:#2}{[S#2]}}
\newcommand{\cataloguescope}[1]{\paragraph{Recorded mathematical question.}#1}
\begin{document}
\section{Explicit bound for symmetric point configurations on the sphere}\label{statement}

\subsection{Definitions and mathematical statement}

\paragraph{Shared notation.}$\mathbb N=\{1,2,\ldots\}$, and $\mathbb Z,\mathbb Q,\mathbb R,\mathbb C$ denote the integers, rationals, reals and complex numbers. Any other conventions are specified locally.


Let $S^2=\{x\in\mathbb R^3:\|x\|_2=1\}$. A nonempty finite subset $X\subset S^2$ is symmetric if a finite group acts transitively on $X$ by isometries of its induced Euclidean metric. A great circle is the intersection of $S^2$ with a plane through the origin; two parallel circles are intersections with two parallel affine planes.
\paragraph{Statement recorded in the supplied source.}
Determine an explicit absolute constant $B$ such that $|X|\le B$ for every symmetric $X$ contained in neither a great circle nor a union of two parallel circles. Is the sharp bound $B=60$, attained by the vertex set of a truncated icosahedron centered at the origin and scaled onto $S^2$? \sref{statement}{2}
\subsection{Short English statement}
Give an explicit size bound for symmetric spherical point sets outside the exceptional circle configurations and decide whether the soccer-ball vertices are maximal.
\subsection{Sources}
\begin{description}
\item[\hypertarget{source:statement:1}{S1}] ULAM.AI and the UnsolvedMath Contributors, \emph{UnsolvedMath: A Curated Collection of Open Mathematics Problems}, catalogue entry. CC BY 4.0. \url{https://huggingface.co/datasets/ulamai/UnsolvedMath}.
\item[\hypertarget{source:statement:2}{S2}] Itai Benjamini, \emph{Points in Equilibrium}, May 2015; original author manuscript preserved in Arquivo.pt capture of 31 December 2020, Unnumbered sphere addendum following Question 0.1, p. 1. \url{https://arquivo.pt/wayback/20201231041612id_/https://www.wisdom.weizmann.ac.il/~itai/equilibrium.pdf}.
\end{description}
\label{end:statement}
\end{document}
